\section{结论}

本文围绕水泥窑尾四电场电除尘器的协同优化，完成了从数据审计、机理建模、工况识别到风险优化的闭环分析，主要结论如下。

\begin{enumerate}
  \item 10,080 个连续分钟中，出口缺失 50 行、54.75\% 的有效读数封顶于 $50\mgNm$。删失 MLE 得到运行点 $50.036\pm0.313\mgNm$，而普通线性解释度仅 0.000815，故附件不能直接识别操作量的排放净效应。
  \item D--A 穿透、温度修正、分场贡献、积灰和再飞扬构成可解释数字孪生。分场权重由中位电压平方归一化得到，温度参考中心由附件均值确定为 $126.008^\circ\mathrm C$，钟形尺度在代码与论文中统一为 $25^\circ\mathrm C$；时间正则化 GMM 得到 8 类工况。能耗代理全样本 $R^2=0.9979$，逐日留一平均 $R^2=0.9866$、RMSE $5.96\kW$。
  \item $10\mgNm$ 下八类工况均通过 95\% 机会约束，总功率为 $1961.4$--$2151.9\kW$。低负荷优先调节 $U_1,U_3$；高负荷四场边际收益接近，应先满足尘负荷与错峰振打安全，再按实时收益协同分配电压。
  \item $5\mgNm$ 在历史边界内有 R1、R7、R8 不可行。电压上界放宽 3\% 的条件情景下，全周期电耗增幅 5.32\%；R8 增幅 5.10\%，95\% 区间 $[4.63\%,7.31\%]$，增量由四场协同承担，其中 $U_1,U_2$ 合计约 65.88\%。
\end{enumerate}

最终应形成“\textbf{低量程监测—负荷前馈—工况滞回—场间分配—振打错峰—风险复核}”的分层控制。可由数据识别的给出精确数值，依赖机理外推的同时给出区间和上线标定条件。
